%% problem_id: S047-AIF-3607-thm-11.2
%% source_id: S047 (NUMDAM; Annales de l'Institut Fourier)
%% source_article: Nicolas Dutertre, Vincent Grandjean, "Equi-singularity of real families and
%%   Lipschitz-Killing curvature densities at infinity", Ann. Inst. Fourier 74 (2024), no. 6,
%%   2379--2423.
%% source_url: https://www.numdam.org/item/10.5802/aif.3607.pdf
%% representation: PDF; statement and proof transcribed by RENDERING the pages (pdftoppm -r 150 -png)
%%   and READING THE IMAGES, then checking every display against them.
%% statement_locator: printed p. 2407 (= PDF p. 30), THEOREM 11.2
%% proof_locator: printed pp. 2407--2409 (= PDF pp. 30--32), from "Proof." through the closing square
%% difficulty_level: H2
%% LEVEL REASONING UNDER THE (NEW) H3 CRITERION -- the decisive step is SPLIT, so H2:
%%   * existence of the open dense subsets is CITED: Omega^k_c := cap_j Omega^{k-s}_{S_j}, where each
%%     Omega^{k-s}_{S_j} comes from Lemma 11.1; the strata come from Claim 11.3, whose own proof is two
%%     sentences ("Following [30, 31], it can be definably stratified").
%%   * carried HERE: the tangent-space computation that shows the specific planes work -- the key fact
%%     K_R + P_1 = R^{n-1}, then N_R := K_R^perp cap R, the projection to R^s restricted to N_R is an
%%     isomorphism, hence R cap (P x R^s) projects surjectively onto R^s; plus the degenerate case
%%     dim W - s = n - k.
%%   Because the harder half (existence) is cited and the half done here is a verification/computation,
%%   my band puts it at H2 under "归约 + 精确记账 / 构造 + 一致性验证".  It is deliberately NOT graded
%%   H3 even though the authors call it "the first main new result of the paper": that is exactly the
%%   over-grading the independent blind check caught twice on R1.
%% status: complete (statement + author proof verbatim + reason + locators)
%% ---------------------------------------------------------------------------
\documentclass[11pt]{article}
\usepackage{amsmath,amssymb,amsthm}
\newtheorem{theorem}{Theorem}
\newtheorem{claim}{Claim}
\begin{document}

\section*{Theorem 11.2}

\textbf{Source.} N.~Dutertre, V.~Grandjean, \emph{Equi-singularity of real families and
Lipschitz--Killing curvature densities at infinity}, Ann.\ Inst.\ Fourier \textbf{74} (2024), no.~6;
printed page 2407.  (The authors introduce it with: \emph{``The first main new result of the paper is
the following:''})

\medskip\noindent
\textbf{Setting.} $\Pi(k) := \{\,P\times\mathbb{R}^s\in\mathbf G(k,n+s) : P\in\mathbf G(k-s,n)\,\}$
is a non-singular projective sub-variety of $\mathbf G(k,n+s)$ isomorphic to $\mathbf G(k-s,n)$.

\begin{theorem}
Let $\mathbf c\notin K(\varphi)$. For every $k\geqslant n+s-\dim W$, there exists a definable open dense
subset $\Omega_{\mathbf c}^k$ of $\Pi(k)$ consisting of $k$-vector planes $P\times\mathbb{R}^s$ such
that the value $\mathbf c$ does not lie in $K(\varphi|_{P\times\mathbb{R}^s})$.
\end{theorem}

\subsection*{Proof (authors' text, printed pp.~2407--2409)}

Observe that for any vector $p$-plane $P$ of $\mathbb{R}^n$ with $p\geqslant 1$, we have
\[
\beta^{-1}(P) = {]0,+\infty[}\times\mathbf S(P)\subset\overline M .
\]
Let us define the ``$p$-plane'' induced from $P$ in $\overline M$
\[
P^{+} := [0,+\infty[{}\times\mathbf S(P)\subset\overline M .
\]
Let $\mathbf c$ be a (regular) value of $\varphi$. The set of accumulation points at infinity of the
value $\mathbf c$ of the mapping $\varphi$ is the closed, definable subset of
$\mathbf S^{n-1}\times\mathbf 0\subset\mathbf S^{n+s-1}$ defined as
\begin{equation}
W_{\mathbf c}^{\infty} :=
\left\{ \mathbf u\in\mathbf S^{n+s-1} \,;\, \exists(\mathbf w_k)_{k\in\mathbb N}\subset W:\;
\mathbf w_k\to\infty,\; \mathbf y_k\to\mathbf c,\; \frac{\mathbf w_k}{|\mathbf w_k|}\to\mathbf u \right\}.
\tag{11.1}
\end{equation}
We write $W^{\infty} = A_{\mathbf c}^{\infty}\times\mathbf 0$. We recall that
$Z = \beta_s^{-1}(W)$ and that $Z^{\infty} = \mathrm{clos}(Z)\cap\overline{\mathbf M}_s^{\infty}$.
Observe that $Z_{\mathbf c}^{\infty} := Z^{\infty}\cap\mathbf 0\times\mathbf S^{n-1}\times\mathbf c
= \mathbf 0\times A_{\mathbf c}^{\infty}\times\mathbf c$.

We recall that $\mathbf z = (r,\mathbf u,\mathbf y)\in\overline{\mathbf M}_s$, and
$\mathbf r_Z(\mathbf z) = r$ for $\mathbf z\in Z$. We need the following intermediary result:

\begin{claim}
Let $\mathbf c\in\mathbb{R}^s\setminus K(\varphi)$. There exists a definable stratification of
$Z_{\mathbf c}^{\infty}$ such that for any point $\mathbf u$ of $A_{\mathbf c}^{\infty}$ and for any
sequence $(\mathbf z_k)_{k\in\mathbb N}$ of $Z$ such that (i) $r_k\to 0$, (ii) $\mathbf y_k\to\mathbf c$,
(iii) $\mathbf u_k\to\mathbf u$, and (iv) $T_{\mathbf z_k}\mathbf r_Z\to R$, we have
\[
T_{\mathbf u}S\subset R,
\]
where $S$ is the stratum of $Z_{\mathbf c}^{\infty}$ containing $\mathbf 0\times\mathbf u\times\mathbf c$.
\end{claim}

\noindent\emph{Proof of the claim.} --- The function $\mathbf r_Z$ extends continuously and definably to
$0$ on $Z^{\infty}$. Following [30, 31], it can be definably stratified $(\mathbf a_{\mathrm{rel}})$.
We can further ask that $\mathbf 0\times A_{\mathbf c}^{\infty}\times\mathbf c$ is a union of strata.
Any stratum of $\mathbf 0\times A_{\mathbf c}^{\infty}\times\mathbf c$ is of the form
$\mathbf 0\times S\times\mathbf c$, for some sub-manifold $S$ of $\mathbf S^{n-1}$.
\hfill$\square$

Let us write again $W_{\mathbf c}\times\mathbf c := \varphi^{-1}(\mathbf c)$. It is more convenient to
work in a neighbourhood of $\overline{\mathbf M}_s^{\infty}$, more precisely nearby
$Z_{\mathbf c}^{\infty}$.

Let $S_1,\dots,S_l$ be the strata of $A_{\mathbf c}^{\infty}$ obtained from Claim 11.3.

Let $k\leqslant n+s$ be any integer such that $k+\dim W-s\geqslant n$. Let us consider the following
definable open dense subset of $\mathbf G(k-s,n)$
\[
\Omega_{\mathbf c}^k := \bigcap\nolimits_{j=0}^{l}\Omega_{S_j}^{k-s},
\]
where $\Omega_{S_j}^{k-s}$ is the definable open dense subset of $\mathbf G(k-s,n)$ of Lemma 11.1
corresponding to $S_j$.

Assume that $\dim W-s\geqslant n-k+1$.

Let $P$ be a $(k-s)$-plane of $\Omega_{\mathbf c}^k$. Let $(\mathbf z_m)_m$ be a sequence of
$P^{+}\times\mathbb R^s\setminus Z^{\infty}$ such that
$\mathbf z_m\to\mathbf z_{\infty} = (0,\mathbf u,\mathbf c)\in Z_{\mathbf c}^{\infty}$. We can assume
that
\[
T_{\mathbf z_m}\mathbf r_Z \longrightarrow R, \quad\text{and}\quad
\mathbf u_m\longrightarrow\mathbf u\in(0\times A_{\mathbf c}^{\infty}\cap P^{+})\times\mathbf 0 .
\]
Since $\mathbf u\in P^{+}$, we deduce that $P = \mathbb R\mathbf u\oplus T_{\mathbf u}\mathbf S(P)$.
All the computations which will follow are done in
\[
T_{\mathbf z_{\infty}}\overline{\mathbf M}_s
= \mathbb R\times T_{\mathbf u}\mathbf S^{n-1}\times\mathbb R^s = \mathbb R\times\mathbb R^{n-1}\times\mathbb R^s .
\]
Therefore we write $P$ for $\mathbb R\times T_{\mathbf u}\mathbf S(P) = T_{(0,\mathbf u)}P^{+}
\subset\mathbb R\times\mathbb R^{n-1}$. By hypothesis $R = 0\times R_1\subset 0\times\mathbb R^{n-1}
\times\mathbb R^s$ and the space $R$ projects surjectively onto $\mathbb R^s$ (see Definition 5.2 and
Lemma 7.1). We want to show that $R\cap P\times\mathbb R^s$ projects surjectively onto $\mathbb R^s$.
Let
\[
0\times K_R\times\mathbf 0 := R\cap(\mathbb R\times\mathbb R^{n-1})\times\mathbf 0
= \ker(R\twoheadrightarrow\mathbb R^s).
\]
We can assume that $K_R = \mathbb R^p\times\mathbf 0\subset\mathbb R^{n-1}$ where $p = \dim W-s-1$.
By hypothesis, denoting by $S$ the stratum of $A_{\mathbf c}^{\infty}$ containing $\mathbf u$, we have
\[
T_{\mathbf u}S\subset K_R, \quad\text{and}\quad P+T_{\mathbf u}S = \mathbb R\times\mathbb R^{n-1}.
\]

Since $T_{\mathbf u}S$ is contained in $0\times\mathbb R^{n-1}$, the $k$-plane $P$ is not. Let
$0\times P_1 := P\cap 0\times\mathbb R^{n-1}$. Since we find
\[
K_R + P_1 = \mathbb R^{n-1},
\]
we deduce the following key fact
\[
\mathbf 0\times\mathbb R^{n-1-p}\subset P_1 .
\]
Let $N_R := K_R^{\perp}\cap R$ be the orthogonal complement of $K_R$ in $R$, thus the projection to
$\mathbb R^s$ when restricted to $N_R$ is an isomorphism. Since the space $N_R$ is contained in
$\mathbf 0\times\mathbb R^{n-1-p}\times\mathbb R^s$, we deduce that
\[
N_R = (0\times P_1\times\mathbb R^s)\cap N_R\subset(P\times\mathbb R^s)\cap R .
\]
In other words $R\cap(P\times\mathbb R^s)$ projects surjectively onto $\mathbb R^s$.

Assume $\dim W-s = n-k$. For each generic $k$-plane $P$ of $\Omega_{\mathbf c}^k$, the fibre
$(\varphi|_{P\times\mathbb R^s})^{-1}(\mathbf c)$ is finite. The definition of
$W_{\mathbf c}^{\infty}$ prohibits, when it is not empty, that $\mathbf c$ be a properness value of
$\varphi|_{P\times\mathbb R^s}$, therefore $\mathbf c$ lies in $K(\varphi|_{P\times\mathbb R^s})$ if
$W_{\mathbf c}^{\infty}$ is not empty. In such a case, the arguments used in the previous case show the
existence of a surjective mapping from a space of dimension $s-1$ onto $\mathbb R^s$.
\hfill$\square$

\end{document}
